\subsection{Sprint 3}

A equipe tinha 5 US previstas, que incluíam muitas tasks de certa complexidade, em geral integrações como Firebase, Mapbox, Sockets e outras. Além disso melhorias de infra e devops fizeram parte do escopo. Da minha parte fiquei responsável pelas tarefas de Backend, principalmente DevOps e integrações. Além disso, no final da sprint, tinha a responsabilidade do deploy do Backend.

Mais uma vez a equipe conseguiu concluir todas as tarefas previstas no escopo, apesar do final caótico de sprint, o resultado da entrega foi novamente perfeito. Da minha parte peguei tasks de integração com mapbox e notificação push, tendo uma leve complicação para fazer as notificações push funcionar principalmente no final da sprint, mas em comunicação com os colegas de equipe consegui entregar a tarefa de modo perfeito. Por fim, dado o code freeze, assumi a responsabilidade de fazer o deploy do backend,   tendo que contornar problemas como falta de armazenamento na S2, mirror quebrado e outros, mas conseguindo depois de muita resiliência entregar.

Por se tratar de uma sprint mais complexa, foi natural o surgimento de diversos problemas conforme o desenvolvimento, até porque, vimos na sprint 3 a oportunidade de deixar tudo que era necessário pronto, para poder entrar em uma sprint 4 de modo muito mais tranquilo, focando apenas em entregas extra e de qualidade de vida para o sistema. Sendo assim problemas no deploy, restrições de linter, problemas com notificações push que requeriam o build oficial de um APK para serem validadas, começaram a aparecer conforme o desenvolvimento, de qualquer jeito a equipe conseguiu se virar perfeitamente, contornando todos problemas encontrados. Da minha parte, vi que precisei gastar mais tempo na última semana da sprint para garantir a qualidade das entregas que fiquei responsável, no entanto, por sorte tive um tempo bom sobrando, que permitiu que a entrega saísse completa.

Imprevistos acontecem e devem ser parte do planejamento, tendo em vista que tendem a surgir cada vez mais, ao se aproximar do final da sprint, tendo em vista que começamos a validar com a régua mais alta, quanto antes começarmos esta validação de régua alta, mais probabilidade tem de entregarmos um conteúdo completo ao final da sprint, outra lição foi a importância de uma code freeze cedo, claro que quanto mais time a equipe tiver para desenvolver melhor, no entanto, estou vendo code freezes com uns 2 dias de antecedência da entrega sendo uma prática frequente e entregando benefícios, dado que é possível revisar problemas e implementar peças faltantes com muito mais calma.

Conversando com os Ages 4, muito trabalho poderia ter sido adiantado para esta sprint, no entanto, dado o calendário optaram por não seguir com isso, senão não teríamos nada para fazer na última sprint. Sendo assim sobraram tasks de devops, alguns nuances no backend, como a troca para JWT e muitas tarefas de documentação, planejo participar muito da parte de documentação, cobrindo arquitetura, infra, tecnologias para estudo e outras áreas que eu tiver capacidade de colaborar. Além disso, junto com os Ages 4 do backend, farei questão de garantir que os últimos ajustes deixem o aplicativo funcionando.